\documentclass[]{article}
\usepackage{todonotes}
\usepackage[T1]{fontenc}
\usepackage{lmodern}
\usepackage{amssymb,amsmath}
\usepackage{ifxetex,ifluatex}
\usepackage{fixltx2e} % provides \textsubscript
\usepackage[utf8]{inputenc}
\usepackage{color}
\usepackage[unicode=true]{hyperref}
\hypersetup{breaklinks=true,
            bookmarks=true,
            pdfauthor={Patrik Jansson et al.},
            pdftitle={GRACeFUL D4.2: A Domain Specific Language (DSL) for GRACeFUL Concept Maps},
            colorlinks=true,
            citecolor=blue,
            urlcolor=blue,
            linkcolor=magenta,
            pdfborder={0 0 0}}
\urlstyle{same}  % don't use monospace font for urls
\setlength{\parindent}{0pt}
\setlength{\parskip}{4pt plus 2pt minus 1pt}
\setlength{\emergencystretch}{3em}  % prevent overfull lines

%% Cezar
\usepackage[margin=1.60in]{geometry}
\usepackage[verbose]{wrapfig}
\usepackage{graphicx}
\usepackage{subfig}
\usepackage{rotating}
\usepackage{lscape}
\usepackage{float}
\usepackage{geometry}
\usepackage{framed}
\usepackage{xspace}
\usepackage{acronym}
\usepackage[square,numbers]{natbib}

\definecolor{GRACeFULblue}{rgb}{0.20,0.60,0.86}

\newcommand{\grace}{GRACeFUL\xspace}
\acrodef{GCM}{\grace Concept Map}
\acrodef{GMB}{Group Model Building}
\acrodef{DSL}{Domain Specific Language}
\acrodef{CFP}{Constraint Functional Program}
\acrodef{RAT}{Rapid Assessment Tool}
\acrodef{CRUD}{Climate Resilient Urban Design}
\acrodef{CLD}{Causal Loop Diagram}
\hyphenation{GRACeFUL}

\author{}
\date{}

\begin{document}

\begin{center}
\includegraphics[width=5cm]{../coverpage/GRACeFULlogo.png}

\textcolor{GRACeFULblue}{Global systems Rapid Assessment tools\\
through Constraint FUnctional Languages}

\vspace{1cm}

FETPROACT-1-2014 Grant Nº 640954

\end{center}

\begin{framed}
\begin{center}
\Large
A Domain Specific Language (DSL) for \\
GRACeFUL Concept Maps\\[1ex]

D4.2\\[1ex]

\end{center}
\end{framed}

\vspace{1cm}

\noindent
\begin{tabular}{@{}ll@{}}
  Lead Participant:       & Chalmers (WP leader: P. Jansson)
\\Partners Contributing:  & KU Leuven, UPC
\\Dissemination Level:    & PU
\\Document Version:       & Draft (version 2017-01-26)
\\Date of Submission:     & 2017-01-31
\\Due Date of Delivery:   & 2017-01-31
\end{tabular}

\newpage

\section*{A Domain Specific Language (DSL) for\\GRACeFUL Concept Maps}\label{DSL4GRACeFUL}

Contributions by: Oskar Abrahamsson, Maximilian Algehed, Sólrún
Einarsdóttir, Alex Gerdes, and Patrik Jansson.

\vfill

\setcounter{tocdepth}{2}
\tableofcontents

\vfill

\newpage

\section*{Abstract}\label{abstract}

This second deliverable (D4.2) of work package 4 presents a Domain
Specific Language (DSL) for describing GRACeFUL concept maps.
%
This is a continuation of the initial work described in ``D4.1 Formal
Concept Maps Elements Descriptions'' delivered in project month 6.
%
The full source code of the language implementation is available on
github and installation instructions are included in this deliverable.
%
The implementation in Haskell can be seen as a formal semantics in
terms of types and functions and this means that GRACeFUL has reached
milestone M4.1 ``DSL with formal semantics v1.0 ready''.
%
In addition we include a section comparing different approaches to
modelling some of the formal semantics concepts relevant for GRACeFUL
concept maps: causal loop diagrams, qualitative probabilistic
networks, and difference equations.

\section{Introduction}\label{introduction}

% explain what we mean with an GCM, restricted version: semi-qualitative version
% of CLDs

% Use a DSL as a kind of intermediate representation

% Not easy to translate from a visual diagram directly to CP

% Overview of Graceful, at least the bit of our concern, place the DSL in context

This document reports on second deliverable (D4.2) of work package 4 of the
\grace project. The main task of this work package is to build a \emph{\ac{DSL}}
for \acp{GCM}. A \ac{GCM} is a representation of policy analysis that contains
all the elements of a policy problem definition, such as goals, criteria, and
a description of the system. A \ac{GCM} is developed and manipulated by
stakeholders during \ac{GMB} sessions. Ultimately a \ac{GCM} is translated
(via the \ac{DSL}) to a \ac{CFP}, which is analysed by a constraint solver. The
results of the constraint solver are presented to the stakeholders.

The \ac{DSL} can be regarded as an intermediate layer between the visual
representation of a \acf{GCM} and a corresponding constraint program.
Translating the visual representation directly to a constraint program is
difficult, because it is hard to check if the generated program is correct. A
DSL alleviates this problem and us to validate the correctness of a model. In
addition, a \ac{DSL} improves the scalability by abstracting away from the
constraint solver. In the longer term this will lead to a \ac{DSL} aimed at
building scalable \acp{RAT} for collective policy making in Global Systems.  

% How have the definitions from D4.1 changed

In the previous deliverable~\cite{d4.1} we have formalised the various elements
of \acp{GCM}. Progressive insight learned us that \acp{CLD} are not
adequate to model the systems the project is envisioning, such as for the \ac{CRUD}
case study. Instead of using \aclp{CLD} we now use
\emph{stock-and-flow} diagram to model the system dynamics. Stock-and-flow
diagrams allow for a more detailed (semi-)quantative analysis.

% We have constructed a DSL, link to github

We have implemented the \acl{DSL} in the functional programming language
Haskell~\cite{haskell98}. Haskell is the so-called host language in which the
\ac{DSL} is embedded. Language features such as algebraic datatypes,
higher-order functions, lazy evaluation, and a rich type system makes Haskell
particulary suitable for defined \acp{DSL}. The details of this language are
explained in Section~\ref{sec:gl}. The actual source code implementing the DSL
is freely available in the project repository on
GitHub\footnote{\url{https://github.com/GRACeFUL-project/}}.

% Explain the structure of the rest of the doc

We continue this document by giving the necessary background information, in
Section~\ref{sec:background}, for making the remaining sections more accessible.
We present the details of the \ac{DSL} for \acp{GCM} including some examples in
Section~\ref{sec:gl}. Finally, we give the formal semantics of the \ac{DSL} in
Section~\ref{sec:semantics}.


\section{Background}
\label{sec:background}

We provide some background information about several aspects we use in
later sections, such that they become more accessible.
%
Whenever possible we give links for further reading as well.

\paragraph{\acfp{GCM}}
In a \acl{GMB} session stakeholders perform a policy analysis that
results in the form of a \acl{GCM}.
%
This concept map contains important elements of a policy problem
definition: goals, criteria for assessing the achievement of those
goals, a description of the system involving factors and criteria, and
the competing alternatives.

The term \emph{factors} refers to characteristics of a system that can
take a value, either quantitative or qualitative, that can change over
time (source: the \emph{GRACeFUL Glossary of Terms}).
%
Similarly, \emph{external factors} correspond to \emph{inputs} to the
system.
%
As such, factors are assumed to be associated with measurable
values.
%
As in systems theory, a distinction is made between inputs that are
beyond the control of the system, and which are potentially uncertain
(or even unknown), and \emph{actions} or combinations of actions,
which represent the system theoretical \emph{controls}.

The interaction between the factors, and between the factors and the
criteria, represents a first description of the system.
%
This description cam be modelled as a hybrid stock-and-flow diagram and 
represents the structural understanding of a
system --- the causal structures that produces the observed behaviour.
%
It reveals information about the rates of change of system elements
and the measures of the variables of the system.
%
%A diagram consists of the following elements:
%\begin{description}
%\item[Stocks] A stock represent a part of a system whose value at any
  %given instant in time depends on the systems past behaviour.
%\item[Flows] Flows represent the rate at which the stock is changing
  %at any given instant, they either flow into a stock (causing it to
  %increase) or flow out of a stock (causing it to decrease).
%\item[Converters] Converters either represent parts at the boundary of
  %the system (i.e. parts whose value is not determined by the
  %behaviour of the system itself) or they represent parts of a system
  %whose value can be derived from other parts of the system at any
  %time through some computational procedure.
%\item[Connectors] Much like in causal loop diagrams the connectors of
  %a system show how the parts of a system influence each other.
%%
  %Stocks can only be influenced by flows (i.e. there can be no
  %connector that connects into a stock), flows can be influenced by
  %stocks, other flows, and by converters.
%%
  %Converters either are not influenced at all (i.e. they are at the
  %systems boundary) or are influenced by stocks, flows and other
  %converters.
%\item[Source/Sink] Sources and sinks are stocks that lie outside of
  %the models boundary --- they are used to show that a stock is
  %flowing from a source or into a sink that lies outside of the models
  %boundary.
%\end{description}
%The elements of a policy problem are described in terms of
%stock-and-flow diagram elements, for example \emph{factors} are
%represented by \emph{stocks}.


\paragraph{\acl{DSL}} \citet{fowler} defines a \ac{DSL} as follows:
%
\emph{a computer programming language of limited expressiveness
  focused on a particular domain}.
%
A \ac{DSL} is targeted at a specific class of programming tasks.
%
By restricting scope to a particular domain, one can tailor the
language specifically for that domain.
%
\citet{Gibbons2015} gives a good overview of implementing \acp{DSL}
using functional programming.
%
There are two main approaches to implementing \acp{DSL}:
\begin{description}
\item[standalone] A language with their own specialised syntax and
  parser to translate programs written in the \ac{DSL} into the host
  language.
%
  An advantage is that the syntax can be tailor made for the target
  audience, and does not have to resemble the host language.
%
  However, creating such a \ac{DSL} is labour intensive and it cannot
  easily reuse features, such as variables and conditionals, from the
  host language.
\item[embedded] An embedded language tries to offer a convenient
  syntax and abstraction mechanisms, but is offered as a library
  written in the host language.
%
  This means that all the existing facilities, such as abstractions
  standard libraries, are directly available.
%
  A disadvantage of an embedded \ac{DSL} is that the users may be
  unfamiliar with the host language.
\end{description}
%
We have chosen for the latter approach because it is better suited
to incorporate changes. Using the `embedded' approach we can quickly
add new features and modify existing functionality. We highlight the 
details of our implementation in Section~\ref{sec:gl}.


\paragraph{Constraint programming} Similar to functional and logic
programming, constraint programming is a declarative programming
paradigm, which means that the order of processing is not fixed.
%
A constraint program is formulated in terms of a number of
constraints.
%
Such constraints are different from the common primitives of
imperative programming languages in that they do not specify a step or
sequence of steps to execute, but rather the properties of a solution
to be found.
%
The program that identifies the solutions satisfying these constraints
is called a constraint solver.
%
In general there might be none, one, or many solutions to a particular
constraint problem.

The constraint programming approach is to search for a state in which
a large number of constraints are satisfied at the same time.
%
A problem is typically expressed as a state containing a number of
unknown variables.
%
The constraint solver searches for values for all the variables.
%
For more information, see, e.g., \url{http://constraint.org}.

The various concepts in a \ac{GCM}, such as factors, are translated 
to constraints in a constraint programming language via the \ac{DSL}. 



%At this stage, if we can attribute qualitative values to the various nodes, the
%causal loop diagram describes a system that can make the object of a qualitative
%simulation. The criteria play the roles of constraints: the qualitative
%trajectories that lead outside the criteria can be pruned by the constraint
%programming system.

%The causal loop diagram is enhanced to a GRACeFUL concept map by the
%addition of extra nodes and non-causal links. The extra nodes represent
%actions, goals, alternatives, and stakeholders.

%Like all previous elements, the actions are also obtained from the
%stakeholders. As explained in the previous section, actions are the
%components of alternatives, among which will be found the policy
%solutions we are seeking. Alternatives are, according to the GRACeFUL
%Glossary of Terms, ``action or combination of actions that influence the
%values of criteria''. They correspond to Walker's ``alternative
%policies'', where a \emph{policy} is defined ``loosely'' as ``a set of
%actions taken to solve a problem''. As in the case of goal definition,
%the GRACeFUL system can suggest such ``atomic actions'' that are useful
%to the goals in the context of the problem at hand.

%The role of actions in the GRACeFUL concept maps is to give initial
%values to the factors and criteria they influence. They do not link via
%causal loops to these factors, since actions are not of the requisite
%type (they do not represent values that can increase or decrease).
%However, they will usually have non-causal links, representing
%constraining relationships to the factors.

%In the other cases of extra nodes (goals, alternatives, and
%stakeholders), the links represent primarily ``belongs to''
%relationships: there are links from stakeholders to goals, indicating
%which stakeholders have declared which goals; there are links from goals
%to the criteria by which they are assessed; finally, there are links
%between alternatives and the factors and criteria that they (i.e., their
%constituent atomic actions) influence.


\section{GenericLibrary: a DSL for GRACeFUL Concept Maps}
\label{sec:gl}

\todo{Make nicer}
GenericLibrary, henceforth abbreviated to GL, is a DSL for descriptions
of GRACeFUL concept map components embedded in the Haskell programming language.
The DSL addresses the issue of bridging the gap between constraint programming
and the visualisation layer by providing abstractions for modular constraint programming.
These abstractions are targeted at simplifying the description of GRACeFUL concept maps.

The DSL is divided in to two parts. The first part, \texttt{GCM} \todo{GCM is a monad},
allows the user to describe the interactions of GRACeFUL concept map components
and has facilities for constructing new components from existing ones \todo{Also known as monadic bind}. The second part,
\texttt{CP} \todo{CP is also a monad}, features primitives for constructing constraint programs
which describe the behaviour of an individual components.

\todo{How do we describe the available language constructs?}

\subsection{The language}\todo{Something about Haskell metaprogramming}
\todo{flesh out?}
As noted earlier, the GL language is primarily constructed from two different languages,
GCM and CP.

\todo{Something about how to include a CP component in a GCM one?}
\todo{\texttt{component :: CP a -> GCM ()}, the \texttt{GCM ()} type ensures nothing directly escapes from the scope of \texttt{CP} to \texttt{GCM}}
The GCM language supports constructing interfaces between components,
and has support for connecting the interface of different components.
The core abstraction in GL is that of the \texttt{Port}, a port is an
entity which represents the way two components interact. A component
exposes some information about the system, a pump may present one port
representing the amount of water being pumped through the pump and another
port representing the maximum flow the pump is able to produce. Ports
may be connected to each other through the \texttt{link} function.
\todo{something about the way in which relationships between factors are represented, i.e. \texttt{linkBy}}

The CP language supports reasoning about integer and floating-point arithmetic, boolean expressions,
and arrays. It has constructions like \texttt{createLVar} and \texttt{assert} for reasoning about
the internal behaviour of a GRACeFUL concept map component.

As can be seen there is a clear separation of concerns in GL between the high-level primitives
for constructing complicated components from simpler ones, represented in GCM, and the low level
implementation details with which the CP language is concerned.

\subsection{Expressing GRACeFUL concept map elements in GL}\todo{Relate to d4.1}
Many of the elements of GRACeFUL concept maps identified in \cite{D4.1} can be modeled in
GL using language primitives such as \texttt{linkBy} (for connections),
\texttt{createAction} (for actions), \texttt{assert} (for constraints) etc.
GL being an \textit{embedded} domain specific language featuring a rich set of primitive
operations for reasoning in the target domain allows the programmer to construct
abstractions which capture behaviour which generalises many of the concepts described
in \cite{D4.1}.


\todo{Say something about more the collection of examples available online: }
The collection of examples is available in the \verb+examples/+ directory of the GitHub repository \href{https://github.com/GRACeFUL-project/GenericLibrary}{GenericLibrary}:

% -- A GCM representing a simluation of the swedish energy system
% energySystem :: GCM ()
% SmallExample.hs

\subsection{Example: Runoff flow}
\label{example-runoff-flow}

We show a small GL program which models a rain runoff area, like a
town square, which has been provided with a pump to alleviate possible
flooding issues (this is a common procedure in countries like the
Netherlands).
%
This example is a small part of a larger model used in the CRUD case
study meant to show how GL can be employed to model concrete problems
in CRUD.


\subsubsection{DSL textual input}
\label{example-runoff-flow-dsl-textual-input}

\todo{some explaining text?}

\begin{verbatim}
pump :: Float -> GCM (Port Float, Port Float)
pump ... -- As before

rain :: Float -> GCM (Port Float)
rain ... -- As before

storage :: Float -> GCM (Port Float, Port Float, Port Float)
storage cap = do
  inflow   <- createPort
  outlet   <- createPort
  overflow <- createPort
  component $ do
    currentStored <- createVariable
    inf <- value inflow
    out <- value outlet
    ovf <- value overflow
    sto <- value currentStored
    assert $ sto === inf - out - ovf
    assert $ sto `inRange` (0, lit cap)
    assert $ (ovf .> 0) ==> (sto === lit cap)
    assert $ ovf .>= 0
  return (inflow, outlet, overflow)

example :: GCM ()
example = do
  (inflowP, outflowP) <- pump 5
  (inflowS, outletS, overflowS) <- storage 4
  rainflow <- rain 10

  link inflowP outletS
  link inflowS rainflow

  output "Overflow" overflowS
\end{verbatim}

When we run the above program we get the output

\begin{verbatim}
ghci> runGCM example
{"Overflow" : 0}
\end{verbatim}
\todo{More explanation needed!}

\todo{Say something about actions!}


\section{Installation and software requirements}
\label{install-and-reqs}

%% -- Installation instructions -----------------------------------------------
%% ----------------------------------------------------------------------------

The purpose of this section is to outline the process of installing
the GRACe software and its require dependencies.
%
Section~\ref{install-overview} provides an overview of the software
dependencies required for development using GRACe.
%
Section~\ref{install-docker} provides installation instructions for a
platform-independent package utilising the
\href{https://www.docker.com/}{Docker}
platform.

%% -- Software dependencies ---------------------------------------------------

\subsection{Software dependencies}
\label{install-overview}

The GRACe language is implemented in
\href{https://www.haskell.org/}{Haskell} and uses the solver tools
from the \href{http://www.minizinc.org/}{MiniZinc} software
distribution.
%
Specifically, development and execution of GRACe programs requires the
following software dependencies to be met:

\begin{itemize}
\item The MiniZinc and Gecode solver software.
%
  These can be found in the MiniZinc software bundle.
%
\item A complete Haskell toolchain able to download packages from
  \href{https://hackage.haskell.org/}{Hackage}.
%
  Two alternative tools (Stack and Cabal) are suitable for this
  purpose, and both are provided by
  \href{https://www.haskell.org/platform/}{the Haskell Platform}.
\end{itemize}

Instructions for installing and using these components are available
in the {GRACe} repository on
\href{https://github.com/GRACeFUL-project/%
  GRACe/blob/master/doc/INSTALL.md}{GitHub}.
%
Alternatively, for those who only wish to run the GRACe
examples, we provide instructions in the following section.

%% -- Using Docker ------------------------------------------------------------

\subsection{Installation using Docker}
\label{install-docker}

Download and install the Docker app from \url{https://www.docker.com/products/docker}.
%
Open a terminal (or the \emph{command prompt} under Windows) and
execute
\begin{verbatim}
  docker pull eugraceful/generic-library
  docker run --rm eugraceful/generic-library
\end{verbatim}

This will download and execute the example located at
\href{https://github.com/GRACeFUL-project/GRACe/blob/master/examples/Examples.hs}{examples/Examples.hs}.


% ----------------------------------------------------------------

\section{Semantics}
\label{sec:semantics}

\todo{Add something about semantics of GenericLibrary?}
\todo{Flesh out this section/transition, refer to D4.1 (or section above describing
  differences from D4.1)}
Originally our main focus was writing a DSL to describe causal loop diagrams (CLDs) and
translate them to system dynamics models, and we spent some time considering
formal semantics that could aid in reasoning about CLDs.

We have modelled CLDs within the \verb|GenericLibrary|, the code for this model
can be seen in the file \verb|QualitativeExample.hs|. \\

\subsection{Causal loop diagrams}

A causal loop diagram (CLD) is a directed graph used to display causal
relationships between variables.
%
The vertices represent the variables and the edges represent qualitative
causal relationships, which can be positive or negative.

We would like to define a semantic that aids us in reasoning about CLDs and will
describe two such semantics here.

\paragraph{Notation}
We denote a positive causal relationship between $A$ and $B$ by
$A\xrightarrow{+} B$ and a negative one by $A \xrightarrow{-} B$.
%
Then $A \xrightarrow{+} B$ means that an increase in $A$ causes an
increase in $B$, while a decrease in $A$ causes a decrease in $B$.
%
On the other hand, $A\xrightarrow{-} B$ means that an increase in $A$
causes a decrease in $B$ (and conversely a decrease in $A$ causes an
increase in $B$).
%
We denote the sign of the edge from $A$ to $B$ by $s_{AB}$, so
$s_{AB}= +$ if $A\xrightarrow{+} B$ and $s_{AB}=-$ if
$A\xrightarrow{-} B$.

A vertex $A$ also has a sign $s_A$ that denotes the total influence on $A$, so
$s_A=+$ if there is an increase in $A$, $s_A=-$ if there is a decrease, $s_A=0$
if there is no change and $s_A=?$ if we cannot determine the change in $A$.

\subsection{QPN approach}

One approach to modelling and reasoning about CLDs is by using qualitative
probabilistic networks (QPNs).

A QPN is defined as a directed
acyclic graph $G=(V,E)$ where the vertices, $V$, correspond to
variables and the edges, $E$ to qualitative probabilistic influences.
%
These influences can be positive (+) or negative (-).
%
Later, when computing transitive influences we will also use (?) for
ambiguous and (0) for probabilistic independence

The meaning of signs on edges is defined according to first order
stochastic dominance, as follows:

Let $F_B(\cdot|a_i, x)$ be the cumulative distribution function (CDF) for
$B$ given $A=a_i$. Then $s_{AB}=+$ means that for all possible values
$a_1,a_2$ of $A$ where $a_1\geq a_2$, we must have:
\[F_B(b_0|a_1, x)\leq F_B(b_0|a_2, x),\]
that is,
\[P(B \leq b_0| A = a_1, x)\leq P(B\leq b_0| A = a_2, x)\]
for all possible values $b_0$ of $B$ and any consistent context $x$.
The context $x$ ranges over all possible assignments to the
variables other than $A$ that influence $B$, that are consistent with both
$A=a_1$ and $A=a_2$.
%
The definition of $s_{AB}=-$ is the same but with $a_1\leq a_2$.

In simpler terms, $s_{AB} = +$ means that greater values of $A$ mean
greater values of $B$ are more likely, and $s_{AB}=-$ means that
greater values of $A$ mean smaller values of $B$ are more likely.

These influences are symmetric, that is, if the edge from $X$ to $Y$ is reversed
we must have $s_{XY} = s_{YX}$.
%
Due to this symmetry it is possible to propagate an observed increase
or decrease of one variable around the graph and find whether other
variables are then likely to have increased or decreased.

TODO: cite Wellman etc.

This definition is broad enough to
apply to many different systems and
be applicable to various real world situations.

We found some issues with QPNs that lead us to explore other
approaches.
%
First of all, since QPNs were originally defined for acyclic graphs
and the theory on them relies on acyclicity, they may not be the best
fit to describe CLDs, in which cycles (feedback loops) are an important
feature.
%
Inference on QPNs containing loops is difficult to implement and can
lead to ambiguous results.

Second, the formal semantics of inference on QPNs is difficult to
formalize since it relies heavily on not-so-simple probability theory.
%
(And stakeholder sessions will not produce data on probability
distributions.)

Additionally, as QPNs are defined solely based on qualitative
relationships there is no obvious way to expand them to also describe
quantitative relationships.

Lastly, since all inference in QPNs is probabilistic it leads to
results that may not be as meaningful or concrete as we would like,
such as ``there is a heightened probability that $x$ has increased'',
rather than ``$x$ has increased''.
%
For instance, a variable may decrease even though the cumulative
probabilistic influence on it is positive.

\subsubsection{Difference equation approach}

Inspired by a system of tanks with water flowing from one to another,
and in search of semantics that might also be extended to quantitative
reasoning, we came up with the following approach.

We consider the values of the graph's vertices to be functions of the same
variable, such as a time variable $t$.

If we have a graph with two vertices, $X$ and $Y$, and one edge from $X$
to $Y$, then $s_{XY}=+$ implies that
\[\frac{\partial Y}{\partial t} = G(X(t)),\]
where $G$ is a monotonically increasing function (monotonically decreasing for negative
causality, $s_{XY}=-$).

If the vertex $Y$ has multiple parent vertices $X_1,\ldots,X_n$, then
$\frac{\partial Y}{\partial t}$ depends on all the parent vertices. We can
isolate the effect of a single parent vertex
 $X$ on $Y$ by differentiation.

In general we can then describe the causal relationship from $X$ to
$Y$ as
\[\frac{\partial\left( \frac{\partial Y(t)}{\partial t} \right)}{\partial Y(t)} =
  g(X(t)),\]

where $g$ has a primitive function $G$ such that $G$ is monotonically
increasing if $s_{XY}=+$ and monotonically decreasing if $s_{XY}= -$.

This is somewhat more nuanced than $CLDs$ as they are described above,
where $s_{XY}=+$ implies that an increase in $X$ leads to an increase
in $Y$, and a decrease in $X$ to a decrease in $Y$.
%
Here we may have some threshold value $x_0$ for $X$, where $G(x_0) = 0$,
above which $X$ always causes an increase in $Y$, but an increase in $X$
causes a faster rate of increase in $Y$ and a decrease in $X$ causes a
slowed rate of increase in $Y$, and vice versa.

Note that though $G(X)$ is monotonically increasing, it may not be strictly
increasing, so we could for instance have $G(X) = 0$ for all $X < C$
for some threshold value $C$.

If the vertex $Y$ has parent vetices $X_1,\ldots,X_n$, then we have
\[\frac{\partial Y}{\partial t} = \sum_{i=1}G_i(X_i),\]
where $G_i$ is monotonically increasing if $s_{X_iY}=+$ but monotonically
decreasing if $s_{X_iY}=-$.

In a discrete time system we consider $\Delta(X_t) = X_t - X_{t-1}$
instead of $\frac{\partial X}{\partial t}$, and write
$\Delta(X_t) = G(Y_{t-1})$ instead of
$\frac{\partial X}{\partial t} = G(Y(t))$.
%
In simple cases we may only consider one time step with two values of
$t$, $t_{start}$ and $t_{end}$.

Here we explore how this approach relates to qualitative reasoning,
but it could be extended to quantitative reasoning by solving
appropriate differential equations.

\paragraph{Simple qualitative model}

We consider a qualitative discrete time system where all values of
vertex variables are either +, -, 0, or ? (where ? is an ambiguous value assigned
to a variable whose value cannot be deduced). These values have the partial
ordering $- < 0 < +$, but ? cannot be compared to the other values. In place of
addition and multiplication we have the operations $\oplus$ and $\otimes$, whose
behavior can be seen in the following tables:
\begin{center}
\begin{tabular}{c|cccc}
$\oplus$ & + & - & 0 & ?\\
\hline
  +   & +  & ? & + & ?\\
  -   & ?  & - & - & ?\\
  0   & +  & - & 0 & ?\\
  ?   & ?  & ? & ? & ?\\
\end{tabular}
\quad
\begin{tabular}{c|cccc}
$\otimes$ & + & - & 0 & ?\\
\hline
  +   & +  & - & 0 & ?\\
  -   & -  & + & 0 & ?\\
  0   & 0  & 0 & 0 & 0\\
  ?   & ?  & ? & 0 & ?\\
\end{tabular}
\end{center}
The only
strictly increasing function in this system is the identity function $id(x) =
x$, and the only strictly decreasing function is the negation function $neg(x) =-\otimes x$.

For simplicity we consider the case where all initial values are
set to zero and $G_e(0)=0$, for all edges $e$. We only consider edge functions
$G_e$ where $G_e(?) = ?$, since we shouldn't be able to make unambiguous
deductions based on ambiguous values.
%

This is convenient for qualitative reasoning since then we are only
concerned with increases and decreases rather than numerical values.
%
The value of variable $X$ at time $t$, which we denote by $X_t$, then tells us whether there has been a net
increase or decrease in $X$.

Consider a graph with three vertices, $Z$ and its two parents
$X$ and $Y$, $X\xrightarrow{s_{XZ}} Z$ and $Y\xrightarrow{s_{YZ}} Z$.
%
Then we have
\[\Delta(Z_t) = G_{XZ}(X_{t-1}) \oplus G_{YZ}(Y_{t-1}),\]
where $G_{XZ}$ and $G_{YZ}$ are monotonically increasing or decreasing in
accordance with $s_{XZ}$ and $s_{YZ}$.
%
If we only allow strictly increasing/decreasing functions we then have
\[\Delta(Z_t) = (s_{XZ}\otimes X_{t-1})\oplus (s_{YZ}\otimes Y_{t-1}).\]

Consider a graph with three vertices $A$, $B$ and $C$ and two edges,
$A\xrightarrow{s_{AB}} B$ and $B\xrightarrow{s_{BC}} C$. Then we have
\begin{align*}
\Delta(C_t) &= G_{BC}(B_{t-1})\\
&= G_{BC}(B_{t-2} \oplus \Delta(B_{t-1}))\\
&= G_{BC}(B_{t-2}  \oplus G_{AB}(A_{t-2}))\\
\end{align*}

If $G_{BC}$ is linear, as is the case when we restrict the available functions
to the strictly increasing/decreasing $id$ and $neg$,
we then have
\[\Delta(C_t) = G_{BC}(B_{t-2})\oplus G_{BC}\circ G_{AB}(A_{t-2}).\]

If we only allow strictly increasing/decreasing functions we then have
\[\Delta(C_t) = (s_{BC}\otimes B_{t-2})\oplus (s_{BC}\otimes s_{AB}\otimes A_{t-2}),\]


\section{Conclusion}

We have designed version 1.0 of a DSL for describing GRACeFUL concept
maps.
%
We have provided a Haskell semantics implementing the DSL and the full
source code is available on github.
%
We have also explored alternative mathematical semantics using
qualitative probabilistic networks and difference equations.

The next actions in work package 4 is
\begin{itemize}
\item implement a middleware for connecting the DSL to the CFP layer,
\item build a testing and verification framework for RATs,
\item assist WP3 in implementing the graphical user interface, and
\item assist WP2 in building up a library of GRACeFUL concept map
  components expressed in the DSL.
\end{itemize}
%
In parallell, the DSL and its implementation will gradually evolve to
express more and more of the requirements extractable from GMB
sessions with stakeholders.



\bibliographystyle{unsrtnat}
\bibliography{d4.2}

\end{document}
